% page 462 of the facsimile (image p-461 of the scan)
% block 157.5 x 246.3 mm ; left margin 8.0 mm
\pgc{-12.700mm}{0.000mm}{
\l{\cel{2}454.}
\l{}
\l{\sou{portar}. (\sou{trans}.)I. Tenar (inter sua manui, inter sua brakii,}
\l{\cel{3}sur sua shultri) ulo pezoza. - II. Prenar kun su, e depo-}
\l{\cel{3}zar en ta o ca loko. EFIS.}
\l{}
\l{\sou{portalo}. I. Quaza vestibulo, multa-kaze ornita per koloni,}
\l{\cel{3}en la enireyo di templo, di kirko, di palaco. - II. Fa-}
\l{\cel{3}sado di kirko en qua esas la pordo precipua. -EF.}
\l{}
\l{\sou{portaveino}. (\sou{anat.}) Veino qua recevas sango de la viceri}
\l{\cel{3}abdominala (ventrala), quan lu distributas aden la he-}
\l{\cel{3}pato. -\cel{2}EF.}
\l{}
\l{\sou{porteo}.\cel{2}Punto til qua la fusilo povas portar la kuglo -}
\l{\cel{3}la kanono, la obuso. -\cel{2}F.}
\l{}
\l{\sou{portfolio}. Sorto di buxo kartona, kunportebla, destinita ye}
\l{\cel{3}recevor paperfolii, desegnuro-folii, dokumenti. - DEFI.}
\l{}
\l{\sou{portiko}. Konstrukturo uzata por la exerci gimnastikala, ek}
\l{\cel{3}fosti qui suportas trabo transversa a qua esas suspendita}
\l{\cel{3}la rigi. - DEFIRS.}
\l{}
\l{\sou{portreto}. Reprezenturo de persono per pikto o desegno sur}
\l{\cel{3}moneto-peco, medalio, telo, e c. - DEFR.}
\l{}
\l{\sou{portuo}. I. Konkavajo di maro aden tero - naturala od artifi-}
\l{\cel{3}cala, ube la navi povas shirmar su. - II. (\sou{metaf}.)Shirm-}
\l{\cel{3}eyo. - III. Loko ube onu embarkas o desembarkas la pasa-}
\l{\cel{3}jeri e la vari transportita per aviacili (aeroplani, heli-}
\l{\cel{3}kopteri, e c.) o navi. - EFIRS.}
\l{}
\l{\sou{portulako}. (\sou{bot.}) Supp-planto di qua la folii esas dika.}
\l{\cel{3}- L.\cel{1}\sou{portulaca}. DEIRL.}
\l{}
\l{\sou{pos}. En epoko, instanto, qua sequas la evento quan onu men-}
\l{\cel{3}cionas. - FL.}
\l{}
\l{\sou{posedar}. (\sou{trans.}) Havar la fakultato juar (ulo od ulu), mem}
\l{\cel{3}sen proprietar lu. - EFIS.}
\l{}
\l{\sou{posho}. Mikra sako de telo, de stofo, sutita en la dikajo di}
\l{\cel{3}vesto, por pozar aden lu to quon onu volas portar kun su.}
\l{}
\l{}
\l{\sou{posibla}. Qua povas eventar. - EFIS.}
\l{}
\l{\sou{posto}. Transporto nacionala di la letri, imprimuri, dokumenti,}
\l{\cel{3}mikra paki. - DEFIRS.}
\l{}
\l{posteno. I. (\sou{milit-arto}). Loko, spaco, okupata da trupo de mi-}
\l{\cel{3}litisti por operaco militala. - II. Plaso asignita a komba-}
\l{\cel{3}tanti. - DEFIRS.}
\l{}
\l{posterno.. (\sou{milit-arto}) Galerio vultoza, subsula, en fortifi-}
\l{\cel{3}kajo por komunikar kun la parto extera. - DEFIS.}
}
